\section{Related Work}
\label{sec:related_work}

The development of terminal-centric AI coding agents draws upon a rich and rapidly evolving body of work spanning code intelligence, autonomous software engineering, and interactive agent design. In this section, we survey the key research streams that inform our system design and contextualize \name within the broader landscape.

\subsection{Code Generation and Code LLMs}
The application of LLMs to programming has progressed from function-level generation, benchmarked by HumanEval~\cite{chen2021evaluatinglargelanguagemodels} and MBPP~\cite{austin2021programsynthesislargelanguage}, through class-level tasks such as ClassEval~\cite{classeval}, to repository-level challenges that demand cross-file planning and multi-step reasoning~\cite{code_survey}. Specialized code LLMs including DeepSeek-Coder~\cite{deepseek_coder}, StarCoder~\cite{starcoder}, CodeLlama~\cite{codellama}, and CodeT5+~\cite{wang2023codet5p} have driven this progression, supported by unified toolkits such as CodeTF~\cite{bui2023codetf} that standardize training and inference across models, while Retrieval-Augmented Generation~\cite{lewis2020rag} and reinforcement learning address the gap between isolated generation and realistic repository-scale tasks~\cite{code_survey}. \name builds on this foundation by embedding code LLMs within an agentic loop that provides the navigation, editing, and execution capabilities necessary to apply generated code in real repositories.

\subsection{Autonomous Issue Resolution}
SWE-bench~\cite{jimenez2024swebenchlanguagemodelsresolve} defined the task of autonomous issue resolution, catalyzing a research frontier that spans single-agent, multi-agent, and workflow-based approaches~\cite{issue_resolution_survey}. Single-agent frameworks such as SWE-Agent~\cite{yang2024sweagentagentcomputerinterfacesenable} pioneered autonomous file navigation and code editing, while AutoCodeRover~\cite{zhang2024autocoderoverautonomousprogramimprovement} and HyperAgent~\cite{phan2024hyperagent} extended this to iterative refinement and generalist task solving. Multi-agent systems distribute the problem across specialized roles: MAGIS~\cite{tao2024magisllmbasedmultiagent} assigns role-playing collaboration, CodeR~\cite{chen2024coderissueresolvingmultiagent} introduces task-graph execution, and platforms such as OpenHands~\cite{openhands} orchestrate heterogeneous agents through ensemble selection~\cite{issue_resolution_survey}. Workflow-based approaches like Agentless~\cite{xia2024agentless} enforce structured pipelines (localization, repair, validation) for improved reproducibility.

Beyond architectural choices, the community has explored both training-based and inference-time methods to improve agent capabilities. Supervised fine-tuning with curriculum learning~\cite{curriculum_learning} and synthetic data~\cite{pham2025swe}, combined with RL algorithms leveraging process-oriented rewards~\cite{process_rewards}, equip models with stronger issue-resolution skills. At inference time, Monte Carlo Tree Search~\cite{mcts_planning} enables flexible backtracking over repair trajectories, while parallel exploration strategies such as CodeMonkeys~\cite{codemonkeys} maximize solution coverage~\cite{issue_resolution_survey}. \name draws on these advances, combining single-agent autonomy with structured subagent delegation and workflow-based safety enforcement. Across these paradigms, agents increasingly rely on code itself as their primary medium of reasoning and action.

\subsection{Code as a Core Medium for Generalist Agents}
Recent work has highlighted a paradigm shift from pure natural language reasoning to code-driven agent interactions~\cite{code_survey}. Using code as a universal medium gives agents precise tool invocation, reproducible state management, and composable action primitives.

\paragraph{Interaction protocols.} Standardized tool-use patterns such as ReAct~\cite{yao2023react} and ReWOO~\cite{xu2023rewoo} enable precise tool invocation and state management. The Model Context Protocol (MCP)~\cite{anthropic2024mcp} introduces structured message formats for reliable multi-turn tool orchestration, while multi-agent coordination schemes such as Agent-to-Agent (A2A) enable direct inter-agent communication~\cite{code_survey}.

\paragraph{Thinking and acting in code.} \emph{Reasoning methods} such as Program-Aided Language Models (PAL)~\cite{gao2023pal}, Program-of-Thoughts~\cite{code_survey}, and Chain-of-Code~\cite{code_survey} enable LLMs to generate and execute code for structured reasoning. \emph{Action execution} frameworks translate plans into runnable code: CodeAct~\cite{wang2024codeact} enables interactive operations through executable Python, TaskWeaver~\cite{qiao2023taskweaver} converts requests into plugin-based function calls, and CodeAgents~\cite{codeagents} provides additional orchestration patterns. \emph{Domain applications} extend this paradigm beyond software engineering, from healthcare (EHRAgent~\cite{ehragent}) to robotic control (Code as Policies~\cite{cap}).

\paragraph{Memory with code.} Code-based storage strategies have proven effective for managing LLM context constraints. Voyager~\cite{voyager} stores validated skills as executable code for later reuse, while Reasoning Bank~\cite{reasoning_bank} enables rule-based learning from repair trajectories. MemGPT~\cite{packer2023memgpt} and ExpeRepair~\cite{experepair} introduce hierarchical and dual-memory architectures for context management, both discussed further in the context engineering subsection below.

By operating within terminal environments and embracing execution feedback through shell commands, OpenDev directly aligns with the ``Acting in Code'' philosophy, leveraging the shell as a universal interpreter to orchestrate complex, multi-step workflows.

\subsection{Agentic Software Engineering Workflows}
Agentic Software Engineering formalizes the workflows through which human engineers and LLMs collaborate on software tasks. Hassan et al.~\cite{sase_survey} provide a comprehensive research roadmap that identifies foundational pillars including agent orchestration, environment design, and lifecycle management.

\paragraph{Iterative and prompt-driven workflows.} The Plan--Do--Assess--Review (PDAR) loop formalizes a single task's lifecycle: planning through Product Requirement Prompts (PRPs), implementation by a dev-agent, self-assessment, and human review~\cite{sase_survey}. CLI toolkits such as SuperClaude template these loops for consistency and convenience, but stop short of a team-level methodology.

\paragraph{Spec-driven development.} A growing family of frameworks treats structured specifications, rather than source code, as the primary source of truth for AI-assisted development. GitHub's Spec Kit~\cite{github2025speckit} introduces a four-phase workflow (Specify, Plan, Tasks, Implement) where a formal \texttt{spec.md} and optional \texttt{constitution.md} govern every agent-driven change, ensuring that AI-generated code aligns with explicitly declared intent and architectural constraints. OpenSpec~\cite{fission2025openspec} takes a complementary brownfield-first approach, designed for evolving existing codebases rather than greenfield projects: each proposed change generates spec deltas (ADDED/MODIFIED/REMOVED requirements in GIVEN/WHEN/THEN format) against a persistent specification baseline, making modifications auditable before any code is written. Both frameworks are agent-agnostic, integrating with 17+ coding assistants via slash commands and filesystem conventions rather than tool-specific APIs.

\paragraph{Multi-agent and agile-inspired frameworks.} BMAD~\cite{blackington2025bmad} takes the team metaphor further by organizing agents into agile roles (Product Owner, Architect, Developer, Scrum Master, Tester), with PRDs and story files enabling parallel execution across isolated Git branches~\cite{sase_survey}. A key architectural decision is work sharding: the Scrum Master agent decomposes tasks into self-contained story files, each carrying exactly the context a Developer agent needs, which addresses context-window limitations by design rather than through compression. A crucial delineation in the broader SASE roadmap proposes separating workspaces into an Agent Command Environment (ACE) for human supervision and orchestration, and an Agent Execution Environment (AEE) for scalable agent operations~\cite{sase_survey}.

\paragraph{Mentorship and lifecycle management.} The SASE roadmap~\cite{sase_survey} introduces the concept of Mentorship-as-Code, where review feedback becomes version-controlled, testable MentorScript rules, enabling cumulative improvement across tasks. Agent lifecycle management shifts agents from stateless contractors to persistent teammates with memory, observability, and secure execution.

\name embodies the characteristics of an execution-focused terminal hub, bridging iterative development cycles natively via command-line interactivity and treating command outputs as direct signals for autonomous debugging and loop refinement.

\subsection{Agent Tool Systems and Modular Components}
In training-free frameworks, LLMs rely on specialized tools to augment reasoning without fine-tuning. These tools are organized along the repair pipeline: bug reproduction, fault localization, code search, patch generation, validation, and test generation~\cite{issue_resolution_survey}.

Bug reproduction tools like AEGIS~\cite{aegis} provide automated environment setup and reproduction workflows, ensuring consistent execution contexts. Fault localization tools include Spectrum-Based Fault Localization (SBFL) and graph-based methods that construct code dependency graphs~\cite{issue_resolution_survey}. Code search tools range from interactive retrieval using BM25 and AST-based APIs to graph-based understanding via Knowledge Graphs and Language Server Protocols~\cite{microsoft2016lsp, issue_resolution_survey}. Patch generation tools employ robust editing formats (e.g., AutoDiff) and ensemble selection through regression testing~\cite{issue_resolution_survey}. SpecRover~\cite{specrover} uses specification-guided search for high-quality patches. Test generation tools such as Otter~\cite{otter} and Issue2Test~\cite{issue2test} use feedback-driven mechanisms to synthesize failing tests that reproduce reported defects~\cite{issue_resolution_survey}.

\name adopts a registry-based tool architecture with lazy discovery, hierarchical skill templates, and defense-in-depth safety mechanisms, balancing comprehensive capability with token efficiency.

\subsection{Context Engineering for Long-Horizon Agents}
Context management is a fundamental challenge for long-running agent systems. Issue resolution tasks often require long-horizon, multi-turn interaction that raises both API cost and performance degradation from context rot~\cite{issue_resolution_survey}.

\paragraph{Formal foundations and taxonomy.} Mei et al.~\cite{mei2025survey} provide the first comprehensive survey of context engineering (CE) for LLMs, formalizing the context supplied to a model as a structured tuple $C = \mathcal{A}(c_{\text{instr}}, c_{\text{know}}, c_{\text{tools}}, c_{\text{mem}}, c_{\text{state}}, c_{\text{query}})$, where $\mathcal{A}$ is an assembly function that orchestrates instructional context, external knowledge, tool schemas, memory, execution state, and the user query. The survey organizes the field around three pillars: \emph{context retrieval} (selecting relevant information from external sources), \emph{context processing} (transforming, compressing, or restructuring retrieved content), and \emph{context management} (maintaining coherence and relevance across turns). \name's architecture maps directly onto this taxonomy: its prompt composer assembles $c_{\text{instr}}$ from modular markdown sections, the tool registry manages $c_{\text{tools}}$ with lazy MCP discovery, the memory pipeline maintains $c_{\text{mem}}$, and adaptive compaction handles $c_{\text{state}}$ across long sessions.

\paragraph{Historical and theoretical perspective.} Hua et al.~\cite{hua2025context} situate context engineering within a broader intellectual history, tracing four eras of development: early prompt engineering focused on single-turn instruction, retrieval-augmented generation that introduced external knowledge, tool-augmented agents that expanded the action space, and the current CE 2.0 era that treats context as a first-class engineering concern. Drawing on Dey's formal definition of context from ubiquitous computing and Weiser's vision of calm technology, they articulate three guiding principles: \emph{entropy reduction} (each context element should reduce uncertainty about the desired output), \emph{minimal sufficiency} (include only what is necessary to avoid attention dilution), and \emph{semantic continuity} (context should evolve coherently across turns rather than being reconstructed from scratch). They also identify a narrowness gap in current systems, which overwhelmingly focus on chat history management while neglecting holistic context dimensions such as tool state, environmental signals, and cross-session knowledge. \name's system reminders, which inject event-driven context at attention-critical positions, directly address the semantic continuity principle, while its staged compaction implements entropy reduction by progressively summarizing low-value history.

\paragraph{Context processing and compression techniques.} The survey literature documents substantial quantitative gains from structured context processing~\cite{mei2025survey}. Chain-of-thought variants (tree-of-thought~\cite{yao2023tree}, graph-of-thought~\cite{besta2023graph}) improve multi-step reasoning by structuring intermediate context, while compression methods reduce token costs without proportional quality loss: the In-context Autoencoder~\cite{ge2023context} achieves 4$\times$ compression, and PREMISE reduces prompt length by 87.5\% while preserving task performance. On the retrieval side, Self-RAG~\cite{asai2023self} introduces self-reflective retrieval that adaptively decides when to retrieve and critiques its own outputs, RAPTOR~\cite{sarthi2024raptor} builds recursive tree-structured summaries for hierarchical retrieval, and GraphRAG leverages knowledge graph structures (e.g., HippoRAG~\cite{gutierrez2024hipporag}) to improve retrieval precision for complex queries. These advances reveal a fundamental asymmetry: LLMs are more robust to compressed context during understanding tasks than during generation tasks, suggesting that aggressive compression is most viable for context that informs reasoning rather than context that directly shapes output text. \name's compaction strategy exploits this asymmetry by aggressively summarizing tool outputs and historical turns (understanding context) while preserving system instructions and recent user messages (generation context) verbatim.

\paragraph{Meta-level context optimization.} Ye et al.~\cite{ye2026meta} introduce Meta Context Engineering (MCE), a framework that treats context engineering itself as an optimization problem rather than a manual design task. They formalize the agent's context as a bi-level optimization: the outer loop searches over context configurations (system prompts, tool schemas, memory strategies) while the inner loop evaluates agent performance on downstream tasks. A key insight is that fixed evaluation harnesses introduce systematic bias into agent benchmarks, as the harness itself shapes the context in ways that may advantage certain strategies over others. MCE addresses this by co-evolving the agent's context and evaluation setup using a $(1{+}1)$-ES with crossover, an evolutionary strategy that mutates and recombines context configurations. On SWE-bench Verified, MCE achieves 89.1\% resolve rate compared to 70.7\% for the hand-engineered baseline, while also being 13.6$\times$ faster due to more efficient context utilization. The optimized configurations transfer across tasks, suggesting that meta-learned context strategies capture general principles rather than task-specific artifacts. This line of work implies that hand-designed context engineering, including the strategies employed by \name, could eventually be complemented or replaced by learned optimization of context assembly.

Memory integration empowers agents to transcend isolated problem-solving by accumulating historical context. Approaches range from hierarchical storage separating general knowledge from repository-specific details~\cite{issue_resolution_survey}, to cognitive architectures that partition knowledge into episodic, semantic, and procedural stores~\cite{mei2025survey, zhong2023memorybank, hu2025memory}. Building on the memory primitives discussed above (MemGPT's virtual paging~\cite{packer2023memgpt}, ExpeRepair's dual-memory~\cite{experepair}), more recent work explores population-level memory: agent variant populations~\cite{agent_populations} maintain diverse exploration histories for robust decision-making, while experience banks~\cite{experience_banks} guide search through accumulated knowledge. Current frontiers prioritize distilling transferable reasoning strategies, shifting from storing raw data to abstracting high-level policies from trajectories~\cite{issue_resolution_survey}.

Memory-driven scaling approaches complement the inference-time strategies discussed above by integrating persistent context to reduce redundant exploration, enabling agents to build on past attempts rather than starting fresh. Context compression for long-horizon agents has seen notable advances: ACON~\cite{kang2025acon} optimizes compression policies for extended agent sessions, while Context-Folding~\cite{sun2025scaling} scales agent capabilities through recursive context summarization. At the frontier, Agentic Context Engineering~\cite{zhang2025agentic} enables models to evolve their own contexts through self-improvement loops. Multi-agent communication has matured through a progression of standardized protocols, from early knowledge-sharing languages (KQML, FIPA ACL) to modern interoperability standards including MCP~\cite{anthropic2024mcp} for tool integration, Agent-to-Agent (A2A) for direct inter-agent messaging, and Agent Communication Protocol (ACP) for structured multi-agent coordination~\cite{mei2025survey}. \name addresses these challenges through automatic compaction that preserves critical information while summarizing history, combined with dual-memory architecture and template-based error recovery. Young~\cite{anthropic2025harness} formalizes the notion of an agent \emph{harness} as the runtime framework that coordinates these concerns (tool dispatch, context lifecycle, progress tracking, and clean handoffs between context windows) for agents operating over extended timeframes.

\subsection{Benchmarks for Evaluating Agentic Coding Systems}

Having introduced foundational code generation benchmarks above, we focus here on the evaluation ecosystem around agentic systems. The scope of benchmarks has progressively widened: from function- and class-level generation (APPS~\cite{apps}, CodeContests~\cite{codecontests}), through repository-level issue resolution anchored by SWE-bench~\cite{jimenez2024swebenchlanguagemodelsresolve} and its variants (SWE-bench Verified~\cite{swebench_verified}, Multi-SWE-bench~\cite{multi_swebench}, SWE-bench Multimodal~\cite{yang2024swebenchmultimodalaisystems}, SWE-bench Pro~\cite{Deng2025SWEBenchPC}, SWE-Lancer~\cite{SWELancer}), to environment-grounded interaction tasks such as WebArena~\cite{zhouWebArenaRealisticWeb2024}, OSWorld~\cite{xieOSWorldBenchmarkingMultimodal2024}, and EnvBench~\cite{eliseevaEnvBenchBenchmarkAutomated2025}. Complementary benchmarks target specific dimensions: SWT-Bench~\cite{SWTBench} for testing workflows, FEA-Bench~\cite{fea_bench} for feature implementation, NL2Repo-Bench~\cite{nl2repo_bench} for repository generation, DevEval~\cite{DevEval} for the full development lifecycle, and SWE-EVO~\cite{thai2025swe} for long-horizon software evolution.

Specialized benchmarks evaluate capabilities beyond general code editing. $\tau$-Bench~\cite{taubench} and the Berkeley Function Calling Leaderboard~\cite{patil2025bfcl} assess tool-use and function-calling proficiency. In scientific and ML engineering, ReplicationBench~\cite{ye2025replicationbenchaiagentsreplicate} targets research replication, MLGym~\cite{mlgym} and MLE-Bench~\cite{MLEBench} evaluate ML research and engineering tasks, and CORE-Bench~\cite{COREBench} measures computational reproducibility. For long-horizon terminal interaction, Terminal-Bench~\cite{terminal_bench} demonstrates that frontier agents resolve fewer than 65\% of curated CLI tasks, while LongCLI-Bench~\cite{longcli_bench} finds pass rates below 20\% for multi-category programming tasks spanning development, feature addition, bug fixing, and refactoring.

\subsection{Evaluation Methodology and Best Practices}

Rigorous evaluation of agentic systems requires careful benchmark design beyond simple accuracy metrics. The Agentic Benchmark Checklist~\cite{agenticbenchmarkchecklist, zhuEstablishingBestPractices2025} formalizes best practices including clear task definition, reproducibility guarantees, contamination prevention, and efficiency-aware metrics. LLM-as-Judge approaches such as MT-Bench~\cite{zheng2023judging} enable scalable quality assessment when traditional metrics fail to capture semantic correctness, though data contamination remains a critical concern requiring continuous benchmark renewal~\cite{contamination_survey, contamination_modern}. Efficiency metrics (API costs, inference time, and token consumption~\cite{efficiency_metrics}) and long-task completion measurement~\cite{long_task_measurement, kwaMeasuringAIAbility2025} provide holistic assessment of agent practicality for real-world deployment.

\subsection{Human--Agent Collaboration}
LongCLI-Bench provides compelling evidence that human--agent collaboration significantly improves task completion. Static plan injection, where key plans are provided before execution, outperforms self-correction in both pass rate and efficiency. Dynamic interactive guidance, where agents request human intervention based on their current state, achieves even higher performance. The combined setting yields the best results while reducing the need for human intervention~\cite{longcli_bench}. These findings strongly indicate that rather than exclusively pursuing full autonomy, future systems should develop collaborative workflows that leverage synergy between efficient agent execution and human strategic guidance. \name supports this paradigm through its approval workflows, interactive command execution, and structured feedback loops.

\medskip
\noindent The research landscape surveyed above reveals a clear trajectory: from isolated code generation toward integrated, long-horizon agent systems that must balance capability, safety, and context efficiency. The benchmarks increasingly demand sustained multi-step reasoning in realistic environments, while the methods progressively address the engineering challenges (context management, tool orchestration, memory, and safety) that such sustained operation requires. The following section synthesizes the specific future directions that emerge from the intersection of \name's architectural experience and these broader research trends.
